%\thispagestyle{empty}
Understanding the radiation pattern of antennas is essential for optimising wireless communication systems. This parameter describes how an antenna radiates or receives electromagnetic energy in space, determining key aspects such as coverage area, signal strength, and susceptibility to interference. Knowledge of the radiation pattern allows better placement and orientation of antennas, minimises energy losses, and improves overall system performance.\\
Although electromagnetic simulation (EMS) software can estimate the radiation pattern during the design phase, practical results often differ from simulations due to environmental influences. Reflections, diffractions, and absorptions can significantly distort the idealised pattern. As a result, understanding the real radiation pattern is crucial to ensure reliable performance in practical applications.\\
However, high-precision antenna measurement equipment is usually very expensive, which makes it impractical for educational laboratories. In this context, this work proposes the development of a low-cost experimental setup for measuring antenna radiation patterns in the Industrial, Scientific, and Medical (ISM) band. The system integrates a two-axis positioning mechanism, an Arduino-based control platform, a radio frequency power detector, and a MATLAB interface for data acquisition and visualisation.\\
Experimental validation using a Yagi–Uda antenna showed good agreement between measured and simulated results. These outcomes indicate the viability of the developed setup.


\noindent {\bf Keywords:} Antenna Radiation Pattern, Wireless Communication, ISM Band, Low-Cost Measurement, Antenna Performance.


%\vfill
%\pagebreak
%\mbox{}
%\vfill
%%\pagebreak